\labsection{3}{Production Planning and Manufacturing in Odoo}{sec:lab03}
\labcovers{ch:production-planning}

\begin{labproject}{Turn a production plan into products, a BOM, and a completed Manufacturing Order}
\textbf{Coverage.} Tutorial 3 + Chapter 5.\\
\textbf{Core system.} Odoo Products, Purchase, Inventory, and Manufacturing.\\
\textbf{Goal.} Connect production-planning concepts to the master data and transactions required to manufacture Road Bikes.

\courselabflow{Create products}{Define BOM}{Produce five}{Reconcile stock}



\begin{labmode}
Tutorial 3 asks for a production approach, forecast, and S\&OP plan. Odoo does not replace those planning calculations in this course. Instead, it demonstrates how an approved plan becomes executable master data and manufacturing transactions.
\end{labmode}

\medskip\noindent\textbf{Part A --- Choose the production approach.}

For your Tutorial 3 business venture, justify make-to-stock, make-to-order, or a more customised approach. For West End Bicycles, use the supplied guide's make-to-stock story: build finished bikes before the customer delivery.

\courselabtrace{One BOM scales to five Road Bikes}{Per bike}{1 frame, 2 tyres, 1 brake, 1 gear, 1 chain}{Consume}{5 frames, 10 tyres, 5 each of the other components}{Produce}{Finished Road Bike stock increases by 5}



\textbf{Use the chapter while coding.} Revisit \hyperref[ch:production-planning]{the earlier chapter explanation}; match its operation to the variables and state changes below.

\textbf{Follow the supplied implementation.} Check the product and units before confirming the manufacturing order. Reservations describe planned availability; they are not the same as consumption. Record any shortage and receipt before marking the production task complete.
\courselabtrace{Trace the five-bike manufacturing order}{1. Input / state}{Road Bike BOM: one frame, two tyres, one brake set, one gear set and one chain per bike.}{2. Operation}{For quantity five, require [5, 10, 5, 5, 5] components in that order.}{3. Check}{After actual consumption and completion, reconcile five finished bikes with the component stock moves.}


\medskip\noindent\textbf{Part B --- Create finished products.}

Create Road Bike and, if required by the class, the additional sample finished products. For Road Bike:
\begin{itemize}
  \item Product Type = Goods;
  \item Track Inventory = enabled;
  \item Can be Sold = enabled;
  \item Sales Price = RM 1,000.
\end{itemize}

\medskip\noindent\textbf{Part C --- Create raw materials.}

Create Aluminium Frame, Rubber Tires, Brake Set, Gear Set, and Bicycle Chain as Goods with Track Inventory enabled and Can be Purchased enabled. Use the guide costs as reference values.

\medskip\noindent\textbf{Part D --- Build the Road Bike BOM.}

Use \pathmenu{Manufacturing \textrightarrow Products \textrightarrow Bills of Materials}. Create a BOM for 1 Road Bike with:
\begin{itemize}
  \item 1 Aluminium Frame;
  \item 2 Rubber Tires;
  \item 1 Brake Set;
  \item 1 Gear Set;
  \item 1 Bicycle Chain.
\end{itemize}

Before manufacturing, verify that components have actually been received into stock. The SCM lab explains purchasing in depth; for this lab, use stock created through the supplied purchase/receipt steps or the instructor's prepared database.

\medskip\noindent\textbf{Part E --- Manufacture 5 Road Bikes.}

Create a Manufacturing Order for quantity 5, assign the Production Planner as Responsible, confirm the order, check component availability, and complete production using the Odoo 19 flow described in the supplied guide. Mark the MO Done.

\medskip\noindent\textbf{Part F --- Verify the inventory transformation.}

Before completion, record component on-hand quantities. After completion, verify:
\begin{itemize}
  \item frame stock decreased by 5;
  \item tyre stock decreased by 10;
  \item each other component decreased by 5;
  \item Road Bike stock increased by 5.
\end{itemize}

This is the ERP evidence that the manufacturing event changed the shared inventory data.

\medskip\noindent\textbf{Part G --- Complete the Tutorial 3 planning work.}

Prepare the sales forecast and S\&OP table required by the tutorial, including the selected product mix and the 90--95\% utilisation requirement specified in that tutorial. Keep the planning units consistent.
\end{labproject}

\textbf{Chapter connection.} Revisit \hyperref[ch:production-planning]{the chapter explanation} and map its states to the steps below.
% Reused from chapters/05-production-planning.tex
\begin{labvisual}{Recall the chapter model}
\begin{labgraphic}
\begin{tikzpicture}[x=1cm,y=1cm]
\foreach \i/\lvl/\q/\h/\d in {%
  0/{Sales forecast}/{What will customers want?}/{12 months}/{product family},
  1/{Sales \& Operations Planning}/{How much can we make, and when?}/{monthly}/{product family},
  2/{Demand management}/{How much of each product, each week?}/{weekly}/{single product},
  3/{Detailed scheduling}/{Which product runs on which day?}/{daily}/{one run}}{
  \pgfmathsetmacro{\y}{-\i*1.12}
  \pgfmathsetmacro{\w}{10.6-\i*0.9}
  \pgfmathsetmacro{\tw}{\w-0.55}
  \node[procstep,minimum width=\w cm,minimum height=0.9cm,text width=\tw cm] (p\i) at (6.0,\y)
    {\lvl\ --- \emph{\q}};
  \node[procchip,anchor=west] at (11.6,\y) {\h};
  \node[procnote,anchor=west,text width=2.3cm] at (13.5,\y) {\d};}
\foreach \i in {0,...,2}{\pgfmathtruncatemacro{\j}{\i+1}\draw[procflow] (p\i)--(p\j);}
\node[procgood,minimum width=3.3cm,minimum height=0.8cm] (prod) at (3.5,-4.6) {Production};
\node[procgood,minimum width=3.3cm,minimum height=0.8cm] (mrp)  at (8.5,-4.6) {MRP $\rightarrow$ Purchasing};
\draw[procflow] (p3.south) -- ++(0,-0.28) -| (prod.north);
\draw[procflow] (p3.south) -- ++(0,-0.28) -| (mrp.north);
\node[proccap,anchor=west] at (11.6,0.78) {\bfseries horizon};
\node[proccap,anchor=west] at (13.5,0.78) {\bfseries detail};
\end{tikzpicture}
\end{labgraphic}\end{labvisual}



\begin{labquestion}{1}{Prepare the component stock}
For a five-bike example, receive at least 5 Aluminium Frames, 10 Rubber Tires, 5 Brake Sets, 5 Gear Sets and 5 Bicycle Chains. Record on-hand quantities before manufacturing.
\end{labquestion}
\begin{labanswer}\textbf{Inspect the result.} Each component is available in the production warehouse. A confirmed purchase order without a completed receipt is insufficient.\end{labanswer}

\begin{labquestion}{2}{Create and complete the manufacturing order}
Open Manufacturing, Operations, Manufacturing Orders, then New. Select Road Bike, quantity 5 and the one-bike BOM. Confirm; inspect component quantities; record the actual consumption and produced quantity, then complete the order.
\end{labquestion}
\begin{labanswer}\textbf{Inspect the result.} The order is Done. Compare each component stock delta with the BOM multiplied by five; Road Bike stock increases by five.\end{labanswer}

\begin{labquestion}{3}{Handle a shortage deliberately}
If a component is unavailable, open its product and receipt records to locate the missing quantity. Complete the receipt or correct the plan before recording production.
\end{labquestion}
\begin{labanswer}\textbf{Inspect the result.} Document the shortage and the upstream record that resolved it. Do not change the BOM merely to conceal unavailable material.\end{labanswer}


\begin{verification}
\textbf{Explain and extend your result.} Reconcile one order from BOM demand to stock moves. Introduce a component shortage in a separate exercise record and explain the resulting availability state and corrective receipt.


\begin{itemize}
  \item The Road Bike and its raw materials exist as Goods with Track Inventory enabled.
  \item The Road Bike BOM lists every component with its quantity.
  \item Components were received into stock before the Manufacturing Order was confirmed.
  \item The Manufacturing Order for 5 Road Bikes is Done.
  \item On-hand stock changed as expected: frames down 5, tyres down 10, other components down 5, Road Bikes up 5.
\end{itemize}
\end{verification}
